\section{Scope, limitations, and conclusions}\label{sec:scope}

\subsection{When a continuum certificate decides a pricing or risk question}
Rough volatility produces both the sharpest and the widest brackets in this study. The exact joint error
ledger of Theorem~\ref{thm:master} identifies the reason before any derivative is priced: the
computable residual variance share $P_\mu/\Var(\Lambda_\mu)$ measures how much continuum randomness is
left unexplained by the retained Gaussian coordinates. For $\mathrm{VIX}_T^2$, every forward variance $\xi_T(u)$
integrates the same past Brownian history on $[0,T]$ through a Volterra kernel analytic across the
diagonal $u=v$, so a single Gaussian direction leaves under $\ShareVIXBase\%$ of $\Var(\mathrm{VIX}_T^2)$ in
standard rough-Bergomi regimes ($\ShareVIXStress\%$ under six-month $\eta_{\mathrm{P1}}=3$ stress). Paired with
Theorem~\ref{thm:tilted} or Proposition~\ref{prop:root} and 169-bit Arb ball arithmetic
(Proposition~\ref{prop:arb}), that low residual share delivers machine-verified futures, option, and
implied-volatility brackets in their corresponding units: the widest of the $61$ certified VIX futures
brackets reported here is $\MaxFutW$ points, below one Cboe tick. For spot integrated variance
$\int_0^T v_t\dd t$ in the same model, fresh rough innovations arrive throughout $[0,T]$, leaving up to
$\ShareIVHi\%$ of variance unexplained by daily closes; the resulting brackets serve as rigorous validation
envelopes rather than quoting tools.

This yields a concrete workflow for volatility desks and model validators:
\begin{enumerate}
\item Compute the residual variance share $P_\mu/\Var(\Lambda_\mu)$ for the target contract and
  proposed Gaussian conditioning space.
\item For VIX futures and options (single or mixed rough Bergomi, flat or curved $\xi_0$), evaluate
  \RoughCertTilt{} (Theorem~\ref{thm:tilted}) or \RoughCertCV{} (Proposition~\ref{prop:arb}). Use fast
  closed-form expansions (such as \citealp{BourgeyDeMarcoGobet2023}, \citealp{LiaoAgarwalBourgey2026}, or
  moment-matched lognormal proxies) inside calibration loops only at parameter cells where
  \RoughCert{} has certified their bias to lie below the desk's tick or implied-volatility tolerance; switch to
  the certified one-factor solver or a verified Markovian factor count wherever the certificate flags a
  violation.
\item When auditing a published Monte Carlo or multilevel benchmark, verify both the continuum target
  and the exact discrete-grid target to separate sampling error from forward-window discretisation bias.
\end{enumerate}

\subsection{Methodological boundaries and honest caveats}
Four boundaries delimit the scope of the present results:
\begin{itemize}
\item \emph{Epistemic distinction between Arb-certified and audited-quadrature results.} Every VIX
  futures, call, put, and Black-76 implied-volatility interval in Sections~\ref{sec:vix_futures}--\ref{sec:vix_options}
  encloses its stated target in the archived data, with zero uncertified
  Gaussian tail truncation and explicit outward rounding. By contrast, the paired quadratic quantile
  benchmarks (Section~\ref{sec:validation}) and the spot integrated-variance studies
  (Section~\ref{sec:variance_risk}) carry \lab{exact-identity+quadrature} and \lab{statistical}$(1-a)$
  labels whose underlying kernel integrals use graded Gauss--Legendre refinement checks rather than
  interval arithmetic.
\item \emph{Separation of asymptotic regimes.} The $O(P^{1/2})$ rate of Proposition~\ref{prop:condvar}(iii)
  describes conditional-variance quantile widths in the local density regime; the $O(\varepsilon^4)$ rate
  of Theorem~\ref{thm:tilted}(iii) describes the relative futures width in the forward-kernel dispersion
  $\varepsilon$; and the $O(\varepsilon^2)$ option bound \eqref{eq:tilted_opt_lip} explains why VIX
  options require conditional variance rather than a static tilt. None of these asymptotic rates replaces
  finite-parameter Arb verification.
\item \emph{Single-driver scope.} All theorems here concern subordinated functionals of one Gaussian
  driver (including mixed rough Bergomi with a shared Brownian motion). Joint SPX--VIX derivatives driven
  by two correlated Brownian motions $(W^{(1)},W^{(2)})$ involve both integrated variance and stochastic
  integrals $\int_0^T\sqrt{v_s}\dd W_s^{(1)}$ and require a separate two-field conditional Monte Carlo
  framework.
\item \emph{Model prices versus market calibration.} All numerical benchmarks evaluate declared
  rough-volatility models at published literature parameters; separating model-pricing error from
  market-calibration misfit requires external market-data inputs beyond the scope of this study.
\end{itemize}

\subsection{Conclusions}
We have introduced \RoughCert{}, an exact Gaussian error calculus and machine-verified certification
framework linking continuum functional approximation to deterministic derivative bounds and conditional risk
certificates. Combining the reconciled Mehler error ledger (Theorem~\ref{thm:master}), conditional-variance
square-root bounds (Proposition~\ref{prop:root}), the static \RoughCertTilt{} Girsanov--H\"older theorem with
sharp $O(\varepsilon^4)$ dispersion width (Theorem~\ref{thm:tilted}), and outward-rounded 169-bit Arb
certification via Pfaff ${}_2F_1$ inner products and Minkowski $L^2$ perturbation bounds
(\RoughCertCV{}, Proposition~\ref{prop:arb}) yields sampling-free continuous-time enclosures narrower than
one Cboe tick across pure, curved-$\xi_0$, and mixed rough Bergomi regimes. Head-to-head audits against
2022--2026 benchmarks machine-verify where order-three weak expansions exceed one full exchange tick on
futures and induce systematic downward implied-volatility bias, quantify direct implied-volatility shortfalls at $52$ of the $60$
published LAB strikes, establish minimum Markovian factor counts for the recorded fits, certify the sign of
the functional-quantiser bias, disentangle
continuum prices from $500$-cell right-rectangle multilevel Monte Carlo targets, and certify steep
upward-sloping VIX smiles in mixed rough Bergomi without simulation.
